\section{Verification, reproduction, and limits}
Each claimed result has five separate components:
\begin{enumerate}
\item A total literal binary table with 215 working states, exact start and halt conventions, and a checked native/canonical bijection.
\item A precise source and its command-calculus calibration. The first uses a tag-permuted strong compiler; the second uses the transfer compiler and its composite-guard refinement.
\item Complete arbitrary-natural primitive contracts, with inductive scan/list lemmas and positive binary durations. The domains include unsafe natural formula codes, which is necessary for both bootstrap proofs.
\item An exact finite blank prefix checked against the whole tape, not just a few decoded registers. The first continues to a safe zero-address root; the second stops at a matched interior AX frame and uses a separately proved total suffix.
\item An all-natural soundness/completeness argument from the calibrated boundary, connected to the explicitly inherited numerical-calculus/ZFC theorem.
\end{enumerate}
Finite source, command, graph, and startup checks support their stated finite claims. They do not substitute for unbounded primitive or scheduler proofs. The mathematical proof engine is executable, but not a proof-assistant kernel. A fresh rerun of that engine is not a second independently designed proof algorithm. Independent parsers, native/dense tape replay, a separately implemented source-prefix check, and mathematical review of the credit proof are identified as such in the receipts.

The archive includes portable entry points for default reconstruction/binding checks and optional complete fresh runtime derivations, together with literal tables, source recipes, compressed lossless traces, proof texts, attribution, and SHA-256 manifests. Its README specifies the exact commands and distinguishes checking saved inventories from rederiving the symbolic proofs. Python assertions must be enabled; entry points reject optimized Python rather than silently skipping checks.

\subsection{What remains open at 214}
The investigated 214-state initialized core still fails blank initialization, including with zero-valued spares. Concrete all-start outward-ray certificates exclude those exact tables and starts, not all possible constructions. Larger source re-encoding and helper-allocation searches also found no certified 214-state blank machine. Timeouts and restricted allocation-model failures are not state lower bounds.

The present deliverable satisfies the alternative criterion of materially different 215-state constructions. It does not rename a failed 214-state table, count source nodes as states, supply a nonblank initial tape, or claim that a source-only compiler already gives a smaller binary machine.

\subsection{Provenance}
The repository is \url{https://github.com/VladimirReshetnikov/ProveIt}. The public zero-spare packet and command-family source were introduced in commit \file{cd596116875}; the transfer-pop packet is the one introduced through commits \file{7342ae8c96c} and \file{2c77edead6a}. The new work combines those inputs and proves the distinct bootstrap calibrations. The underlying numerical calculus, dominated-row certificate, arithmetic architecture, command compilers, credit theorem, and runtime induction checker retain their upstream attribution and licenses.

The primary numerical-calculus soundness boundary is reached only after the complete 570,263-step finite prefix. The secondary boundary contains only eight assumed theorem heads, not all ten AX outputs. These qualifications are part of the theorem, not merely implementation notes. Neither the repository's uncompiled endpoint drafts nor a literal-data audit is cited as a new kernel-checked endpoint for these tables.
